% PROPOSTA-DE-RESOLUCAO.tex -- proposta de norma da Escola Politecnica/UFRJ
% para o Projeto de Graduacao, no lugar do Anexo da Resolucao n. 05/2012.
%
% O documento de apoio e o NAO-CONFORMIDADES-POLI.md, na raiz do repositorio:
% ele confere a Resolucao 05/2012 contra o Manual UFRJ 2026, item a item, e e
% de la que sai o que esta proposta guarda e o que ela descarta.
%
% Usa a classe article de proposito, e o mesmo desenho da NORMA_COPPE_2026.tex:
% e um documento normativo SOBRE a classe, nao um Projeto de Graduacao, e
% aplicar a propria classe a ele confundiria.
%
% NAO sai do ufrj.dtx nem do ufrj-poli.dtx: a norma tem ciclo de vida proprio,
% decidido pela Congregacao da Escola, e nao e artefato da classe.

\documentclass[a4paper,11pt]{article}
\usepackage[T1]{fontenc}
\usepackage[utf8]{inputenc}
\usepackage{lmodern}
\usepackage[brazilian]{babel}
\usepackage[margin=2.5cm]{geometry}
\usepackage{titlesec}
\usepackage{enumitem}
\usepackage{microtype}
\usepackage{url}
\usepackage{parskip}
\usepackage{longtable}
\usepackage{booktabs}
\usepackage{array}
\usepackage[hidelinks]{hyperref}

\titleformat{\section}{\large\bfseries}{\thesection}{1ex}{}
\titleformat{\subsection}{\normalsize\bfseries}{\thesubsection}{1ex}{}

% Cada secao normativa declara de que tipo e a diferenca em relacao ao Manual.
\newcommand\tipo[1]{\normalfont\textit{--- #1}}

\title{\bfseries Norma para a Elaboração Gráfica do Projeto de Graduação
       da Escola Politécnica/UFRJ\\[2mm]
       \large Edição 2026 -- Proposta para apreciação da Congregação}
\author{}
\date{Setembro de 2026}

\begin{document}
\maketitle

\noindent\textit{Esta Norma, quando aprovada, substitui o Anexo da} Resolução
n.\textsuperscript{o} 05, de 2012 \textit{--- Estabelece Normas de Elaboração
Gráfica do Projeto de Graduação.}

\hrulefill

\section*{Como ler esta Norma}

Esta Norma \textbf{não descreve como se formata um Projeto de Graduação}. Quem
faz isso é o \textbf{Manual para Elaboração e Normalização de Trabalhos
Acadêmicos} da UFRJ/SiBI, 9.\textsuperscript{a} edição revista (Rio de Janeiro,
2026) --- daqui em diante \textit{Manual UFRJ 2026} ---, que vale integralmente
na Escola Politécnica.

O que esta Norma faz é registrar, ponto a ponto, \textbf{onde a Escola
especializa o Manual}: onde ele deixa uma escolha em aberto e a Escola escolhe,
onde ele pede um dado que só a Escola tem, e onde a Escola acrescenta um
elemento que ele não prevê. Fora do que está escrito aqui, vale o Manual ---
formato do papel, margens, espaçamento, ordem dos elementos pré-textuais,
numeração das folhas, ilustrações, citações e referências.

Cada seção abaixo diz de que tipo é a diferença:

\begin{description}[leftmargin=2.6cm,style=nextline,font=\bfseries]
  \item[Escolha] O Manual admite mais de uma forma; a Escola fixa uma.
  \item[Dado próprio] O Manual pede a informação; o valor é da Escola.
  \item[Acréscimo] Elemento que o Manual não prevê e a Escola exige ou admite,
    onde o Manual não obriga e sem contrariar exemplo dele.
  \item[Reafirmação] Nada é acrescentado; repete-se um ponto do Manual por ser
    fonte frequente de erro.
\end{description}

\textbf{O que esta Norma pode decidir.} Ela \textbf{não contraria o Manual}:
onde o Manual obriga, vale o Manual. Ela decide \textbf{só onde o Manual não
obriga} --- onde ele deixa uma escolha em aberto ou pede um dado que só a
Escola tem ---, e, ao decidir, \textbf{fica com o que os exemplos do Manual
mostram}. As duas folhas do modelo que o SiBI distribui --- a folha de rosto e
a folha adicional, que a 3.1.2.1.2 manda baixar --- valem como \textbf{errata}
sobre os exemplos do Manual. As normas da ABNT valem nas edições que o Manual
cita, e por meio dele (Seção~\ref{norma:citacoes}).

\textbf{Por que uma Norma nova.} O Anexo da Resolução 05/2012 é anterior ao
Manual vigente e o contraria em treze pontos --- entre eles a numeração das
folhas pré-textuais, a ordem dos elementos, a posição da legenda das figuras, a
folha de aprovação e, sobretudo, o formato das referências, que ele baseia na
NB-66, norma da ABNT substituída em 1989. Cada um desses pontos está conferido
no documento \texttt{NAO-CONFORMIDADES-POLI.md} que acompanha esta proposta.
Desta Norma consta apenas o que continua sendo decisão da Escola.

\subsection*{Mapa das diferenças}

\begin{longtable}{@{}p{0.42\linewidth}p{0.28\linewidth}p{0.22\linewidth}@{}}
\toprule
\textbf{Seção desta Norma} & \textbf{Item do Manual} & \textbf{Tipo} \\
\midrule
\endhead
1. Logotipos & Anexo A & Escolha\\
2. Nome da instituição na capa & Anexo A, 3.1.2.1.1 & Escolha\\
3. Cursos de graduação da Escola & 3.1.2.1.1(e) & Dado próprio\\
4. Título conferido & 3.1.2.1.1(e) & Dado próprio\\
5. Identidade institucional em português & 3.1.2.1.1, 3.1.2.1.3 & Escolha\\
6. Folha de resumo & 3.1.2.1.4, Anexos E e F & Reafirmação\\
7. Resumo em português e em inglês & 3.1.2, 3.1.2.1.4, 3.1.2.1.5 & Escolha\\
8. Extensão do resumo & 3.1.2.1.4 & Escolha\\
9. Família da letra & 2.2 & Escolha\\
10. Folha adicional e ficha catalográfica & 3.1.2.1.2 & Reafirmação\\
11. Citações e referências & 4.1, 4.2 & Escolha\\
12. Idioma do trabalho & 2.1 & Escolha\\
\bottomrule
\end{longtable}

\section{Logotipos \tipo{Escolha}}

O Anexo A do Manual põe, no alto da capa, o \textbf{logotipo da UFRJ à
esquerda} e o \textbf{logotipo da unidade à direita}, ambos marcados como
opcionais.

\textbf{A Escola adota os dois}, e o da direita é a \textbf{marca principal da
Escola Politécnica, colorida}, e não a versão em preto e branco que circula em
documentos impressos --- o preto e branco é limitação da impressão, e não a
marca. Os dois logotipos têm a mesma altura e figuram \textbf{somente na capa}:
a folha de rosto não os repete, como mostram o Anexo B do Manual e as duas
folhas do modelo do SiBI.

\section{Nome da instituição na capa \tipo{Escolha}}

A capa nomeia, em três linhas centralizadas e em caixa alta: a
\textbf{Universidade Federal do Rio de Janeiro}, a \textbf{Escola Politécnica} e
o \textbf{Curso}. O nome do Curso segue a forma da Seção~\ref{norma:cursos}.

\section{Cursos de graduação da Escola \tipo{Dado próprio}}
\label{norma:cursos}

O Manual pede o nome do curso a que o trabalho é submetido, e não tem como
saber quais são. São os treze cursos de Engenharia da Escola e o curso de
Nanotecnologia:

\begin{center}
\begin{tabular}{@{}ll@{}}
\toprule
\textbf{Curso} & \textbf{Título conferido} \\
\midrule
Engenharia Ambiental & Engenheiro Ambiental\\
Engenharia Civil & Engenheiro Civil\\
Engenharia de Computação e Informação & Engenheiro de Computação e Informação\\
Engenharia de Controle e Automação & Engenheiro de Controle e Automação\\
Engenharia Elétrica & Engenheiro Elétrico\\
Engenharia Eletrônica e de Computação & Engenheiro Eletrônico e de Computação\\
Engenharia de Materiais & Engenheiro de Materiais\\
Engenharia Mecânica & Engenheiro Mecânico\\
Engenharia Metalúrgica & Engenheiro Metalúrgico\\
Engenharia Naval e Oceânica & Engenheiro Naval\\
Engenharia Nuclear & Engenheiro Nuclear\\
Engenharia de Petróleo & Engenheiro de Petróleo\\
Engenharia de Produção & Engenheiro de Produção\\
Nanotecnologia & \textit{a definir}\\
\bottomrule
\end{tabular}
\end{center}

\textbf{A forma do nome é ``Curso de Engenharia X''}, onde X é o nome do curso,
e ela é a mesma em toda folha em que o curso aparece: na capa
(``CURSO DE ENGENHARIA DE PETRÓLEO''), na frase da natureza (``apresentado ao
Curso de Engenharia de Petróleo''), na ficha catalográfica
(``UFRJ/Escola Politécnica/Curso de Engenharia de Petróleo'') e na referência
do alto do resumo (``Graduação em Engenharia de Petróleo'').

\section{Título conferido \tipo{Dado próprio}}

A 3.1.2.1.1(e) pede o \textbf{objetivo} do trabalho --- o grau pretendido ---, e
ele é o título que o curso confere, na tabela da Seção~\ref{norma:cursos}. A
frase da natureza fecha com ele: ``\ldots\ como requisito parcial à obtenção do
título de Engenheiro de Petróleo''.

\section{Identidade institucional em português \tipo{Escolha}}

A capa, a folha de rosto, a folha adicional e a folha de aprovação são
\textbf{identidade institucional} e ficam \textbf{em português}, mesmo quando o
trabalho é escrito em outro idioma. Acompanham o idioma do trabalho apenas o
título e o subtítulo. É a mesma escolha da COPPE, e pela mesma razão: são
folhas que a Universidade lê.

\section{Folha de resumo \tipo{Reafirmação}}

Cada folha de resumo segue o modelo dos Anexos E e F do Manual: o título da
folha (RESUMO, \textit{ABSTRACT}), a referência do trabalho, o texto do resumo
e as palavras-chave. \textbf{A Escola não acrescenta nada a ela.}

A Resolução 05/2012 acrescentava a frase de abertura (``Resumo do Projeto de
Graduação apresentado à Escola Politécnica/UFRJ\ldots''), o título com o nome
do autor e o mês, o orientador e o curso. Nada disso está no modelo do Manual,
e a referência já traz o autor, o título, a natureza, o curso e o ano. A
implementação de referência continua a oferecer esses elementos como opção,
desligada por padrão, para o caso de a Escola voltar a exigi-los.

\section{Resumo em português e em inglês \tipo{Escolha}}

O Manual prevê dois resumos, nesta ordem (3.1.2): o \textbf{resumo em língua
vernácula} (3.1.2.1.4), que é o português, e o \textbf{resumo em língua
estrangeira} (3.1.2.1.5). A Escola fixa o \textbf{inglês} como a língua
estrangeira, como já fazia. \textbf{O resumo em português vem sempre primeiro},
qualquer que seja o idioma do trabalho.

\section{Extensão do resumo \tipo{Escolha}}

O Manual admite até 500 palavras (3.1.2.1.4). A Escola pede \textbf{até 250},
como na Resolução 05/2012: é uma restrição dentro do limite do Manual, e não o
contraria.

\section{Família da letra \tipo{Escolha}}

O Manual fixa o \textbf{corpo 12} (2.2b) e não nomeia família. A Escola admite
\textbf{Times New Roman ou Arial}, e admite também a fonte padrão da
implementação de referência. O corpo é 12 em todas: o ``Arial 11'' da Resolução
05/2012 está superado pelo Manual.

\section{Folha adicional e ficha catalográfica \tipo{Reafirmação}}

A ficha catalográfica vem na \textbf{folha adicional}, gerada pelo Gerador de
Fichas Catalográficas do SiBI (\url{http://fichacatalografica.sibi.ufrj.br/}),
como manda a 3.1.2.1.2, e não desenhada à mão. Os campos da Coleta CAPES dessa
folha são das teses e dissertações, e não do Projeto de Graduação: a folha do
Projeto de Graduação traz \textbf{apenas a ficha}.

\section{Citações e referências \tipo{Escolha}}
\label{norma:citacoes}

Valem as normas da ABNT \textbf{nas edições que o Manual cita} na sua lista de
referências --- \textbf{NBR 6023:2025} (Referências) e \textbf{NBR 10520:2023}
(Citações) ---, e não ``as edições vigentes''. Onde uma delas diverge do
Manual, vale o Manual.

Qualquer dos dois sistemas de chamada que o Manual prevê --- autor-data ou
numérico (4.1.1.1) --- é admitido, \textbf{a critério do autor}, desde que
usado de forma consistente em todo o trabalho; o numérico não se usa em
trabalho com notas de rodapé (4.1.1.1.1). É a mesma liberdade que a Resolução
05/2012 dava.

\textbf{O que não vale mais é o formato das referências daquela Resolução},
baseado na NB-66: o número entre colchetes precedido do sobrenome, o ``pp.''
antes das páginas, o título entre aspas, o grifo do título do periódico, o
``et al.'' a partir do quarto autor e o endereço eletrônico entre
\texttt{<\,>}. Todos esses pontos passam a seguir a NBR 6023:2025, pelo Manual.

\section{Idioma do trabalho \tipo{Escolha}}

O trabalho é escrito em português. Inglês ou espanhol são admitidos mediante
justificativa ao Curso, o que está de acordo com o art.~57 da Resolução CEPG
n.\textsuperscript{o} 302/2024 e com a 2.1 do Manual. A escolha afeta o corpo
textual, as legendas, os títulos das seções pré-textuais e o título do trabalho
na capa e na folha de rosto; não afeta a identidade institucional (Seção 5).

\section{Implementação de referência}

A Escola distribui, como implementação de referência desta Norma, a classe
\texttt{ufrj} com o estilo \texttt{ufrj-poli}, das versões 5.0 e posteriores do
\textbf{CoppeTeX}, em \url{https://github.com/COPPE-UFRJ/CoppeTeX}. Ela compõe
o Projeto de Graduação segundo o Manual UFRJ 2026 e as decisões acima:

\begin{verbatim}
  \documentclass[grad]{ufrj}
  \usepackage{ufrj-poli}
  ...
  \department{PETROLEO}
\end{verbatim}

O trabalho escrito com a classe \texttt{poli} anterior continua compilando: a
classe \texttt{poli} passa a ser a classe \texttt{ufrj} com este estilo, e
aceita as assinaturas antigas de \verb|\advisor|, \verb|\coadvisor| e
\verb|\examiner|.

\subsection{O que a implementação não faz}

A lombada (NBR 12225), a contagem de palavras do resumo e o que é da via
impressa --- papel, gramatura, encadernação e cor da capa --- ficam por conta
do autor e da Secretaria de Graduação.

\end{document}
